\section{The \instname\ data reduction flow.}
\label{sec:reduction_flow}


The overall data flow of the \instname\ pipeline is displayed in Figure \ref{fig:reduction_cascade}.

 
\begin{figure}[h!]
\centering
\includegraphics[width=0.9\textwidth]{../instrument_tutorial/figures/reduction_cascade.jpg}
\caption{The data reduction cascade of the \instname\ workflow.}
\label{fig:reduction_cascade}
\end{figure}



The reduction cascade is organized in tasks, which represent
well-defined steps in the process. Tasks can be grouped inside
sub-workflows.  Each task runs a recipe; the detailed description of
the algorithms, input, outputs and recipe parameters used in each
recipe are available in the pipeline manual. Here, we present only the
description of most important features.

The \wkfname\  \edps\ workflow is designed to execute the tasks that
deliver the final reduced data cube for each dataset. It can be either
the product of a single exposure, or the combination of multiple
exposures. Only calibrations needed by the selected the scientific
exposures are processed.

It is possible to set \edps\ to perform the data reduction until a
certain step of the reduction chain (e.g. to reduce only standar
stars, or only flat fields).  This is done by specifying the desired
tasks in the field {\bf Select reduction target} of the {\bf Raw Data}
tab.

The reduction steps of the \wkfname\ workflow are listed below. Before
starting the reduction, the parameters of the recipes associated to
each task can be configured by pressing the button

\includegraphics[width=0.8cm]{../common/figures/configure_dataset.jpg}
close to each dataset configuration.
See for more info on the configuration editor \ref{sec:configuration_editor}






\subsection{Flats and Tracing}

The subworkflow {\bf flats and tracing} is responsible to determine
the traces of the various orders on the detector, determine the order
profile, the pixel-to-pixel efficiency variations, and the blaze
function. These steps are carried by the following 3 tasks within the
subworflow:

\subsubsection{Led flats}

The task {\bf led\_flats} runs the recipe {\tt espdr\_led\_ff} to
process led exposures and determine maps of bad pixels.  In a typical
dataset, led flat fields are not present and static calibrations that
map the bad pixels (whose location is very stable) are used instead.


\subsubsection{Order Definition}

The task {\bf order\_definition} runs the recipe {\tt espdr\_orderdef}
to trace the various orders on the detector by using thungsten
lamps. Typically, the default values give already good results. In
some rare occasions, the presence of strong cosmic rays makes the
recipe to fail; there are no recipe parameters that deal with cosmic
rays rejection, therefore cleaning the input files before processing
is the only option if the problem occurs.

\subsubsection{Flat Fielding}

The task {\bf lamp\_flats} runs the recipe espdr\_mflat and generates
calibrations needed to correct for pixel to pixel variation and the
blaze function. Typically, the default values give already good
results. If ``suspicious bumps'' are seen in the masterflat or blaze
function, it could be an indication that the intra-order background
was badly subtracted (very rare). In this case, the recipe parameters
that can be changed to improve the results are:

\begin{verbatim}
background_sw       : Background subtraction activation (on/off). [on]
bkgr_grid_size_x    : Grid size in x used to calculate the background,
                      between: 16 and 10000. [577]
bkgr_grid_size_y    : Grid size in y used to calculate the background,
\end{verbatim}

\subsection{Wavelength calibration}

The subworkflow {\bf wavelength\_calibration} executes two tasks {\bf
  wave\_thar\_fp} {\bf wave\_thar\_thar}, which both run the recipe
{\tt espdr\_wave\_TH\_drift}. They generate the wavelength solution to
be applied to the science data using Thorium-Argon lamps and
Fabry-Perot exposures.

 Typically, the default values give already good results.  In some
 rare occasions, the presence of strong cosmic rays makes the recipe
 to fail; there are no recipe parameters that deal with cosmic rays
 rejection, therefore cleaning the input files before processing is
 the only option if the problem occurs.

In some rare cases, the peaks of the emission lies are not enough
intense and the recipe can miss some identifications. In the
unfortunate event that none of the peaks that should be identified in
both orders are detected, the recipe cannot "connect" the wavelenght
solution between orders and it fails. The error message in the recipe
log in this case is:

\begin{verbatim}
[ ERROR ] espdr\_wave_THAR: espdr_get_all_FP_ll_whole_image failed:
Input data do not match
\end{verbatim}

In some cases, this situation can be avoided by changing the recipe
parameter {\bf fp\_extraction\_ksigma}.

\subsection{Contamination Thar and Contamination FP}

These two tasks determines the amount of mutual contamination between
the two fibres when exposes with Thorium-Argon or Fabry-Perot lamps by
running the recipe {\tt espdr\_cal\_contam}. Typically, only data to
determine the contamination from Thorium-Argon lamps are present.
Default values give already good results.

\subsection{Flux calibration}

\subsection{Other tasks and workflow not associated to the scientific reduction}

\subsubsection{Efficiency}
\subsubsection{Detector properties}

\begin{itemize}

\item {\bf Detector\_monitoring}

\item{\bf Dark} This task runs the recipe {\tt
  espdr\_mdark} to determine hot pixel mask. In a typical scientific
  dataset, raw darks are not present and static calibrations that map
  the hot pixels (whose location is very stable) are used instead; the
  reduction of dark frames is mainly used for instrument monitoring.
\end{itemize}
